\section{Introduction}\label{sec:intro}

\subsection{The Euclid--Mullin sequence}

The \emph{Euclid--Mullin sequence} is defined by $a_0 = 2$ and
\[
  a_{n+1} = \text{least prime factor of } P(n) + 1,
  \qquad P(n) = a_0 \cdot a_1 \cdots a_n.
\]
The first several terms are $2, 3, 7, 43, 13, 53, 5, 6221671, \ldots$
(OEIS~A000945). The construction mimics Euclid's proof of the infinitude of
primes: each $P(n)+1$ is coprime to all previously seen primes, so
$a_{n+1}$ is always a \emph{new} prime. Mullin~\cite{Mullin63} asked
whether every prime eventually appears in this sequence.

\begin{openproblem}[Mullin's Conjecture, 1963]
For every prime $q$, there exists $n$ such that $a_n = q$.
\end{openproblem}

Despite substantial computational evidence---Guy and Nowakowski~\cite{GuyNow75}
computed 51 terms, and subsequent work has extended this further---Mullin's
Conjecture remains open. The prime $41$ has not appeared in the first several
hundred computed terms, and Cox and van der Poorten~\cite{CoxVdP} raised the
possibility that it might never appear.

\subsection{The function field analog}

Over the polynomial ring $\Fpt$ (with $p$ a prime number), monic irreducible
polynomials play the role of primes. We define the \emph{function field
Euclid--Mullin sequence} (FF-EM) as follows.

\begin{definition}[FF-EM Sequence]\label{def:ffem}
\leancheck{FFEMData}
Fix a prime $p$. The FF-EM data consists of:
\begin{itemize}
  \item $\ffseq(0) = t$ (the ``smallest'' monic irreducible),
  \item $\ffprod(0) = t$ (the running product),
  \item $\ffseq(n+1)$ is a monic irreducible factor of $\ffprod(n)+1$ of
    minimal degree,
  \item $\ffprod(n+1) = \ffprod(n) \cdot \ffseq(n+1)$.
\end{itemize}
\end{definition}

\begin{openproblem}[Function Field Mullin Conjecture]
\leancheck{FFMullinConjecture}
For every monic irreducible polynomial $Q \in \Fpt$, there exists $n$ such
that $\ffseq(n) = Q$.
\end{openproblem}

Two structural differences from the integer setting are immediate.
First, among irreducible polynomials of the same degree, there is no
canonical total ordering; the minimality condition in Definition~\ref{def:ffem}
is on \emph{degree} only, so the successor $\ffseq(n+1)$ may be chosen
among several competitors. Second, the unit group
$(\Fpt/(Q))^\times \cong \GF{p}{\deg Q}^\times$ is cyclic of order
$p^{\deg Q} - 1$, providing a richer group-theoretic structure than the
integer analog $(\ZZ/q\ZZ)^\times$.

\subsection{Why study the function field case?}

The function field setting enjoys a decisive analytic advantage over
$\ZZ$: the Weil bound for character sums over $\Fp$-curves is a
\emph{theorem} (Weil, 1948~\cite{Weil48}), not a conjecture. This
makes population equidistribution of monic irreducibles across residue
classes unconditional, in contrast to the integer case where the
analogous result requires the Weighted Prime Number Theorem in
arithmetic progressions (a consequence of the Generalized Riemann
Hypothesis for quantitative error terms, or of Siegel--Walfisz for
qualitative versions).

By working over $\Fpt$, we can isolate the \emph{orbit-specificity
barrier}---the gap between knowing that irreducibles are equidistributed
across residue classes (population statistics) and knowing that the
\emph{specific} walk $\ffprod(n) \bmod Q$ visits the residue class $-1$
(orbit dynamics). Our results demonstrate that this barrier is the
\emph{sole remaining obstacle} to FF-MC, and that it is structurally
identical to the integer case.

\subsection{Main results}

We prove four main theorems, all unconditional (zero open hypotheses in the
formal verification).

\begin{theorem}[Almost-All Generalized Mixed MC --- Theorem~D]\label{thm:main-sieve}
\leancheck{ff\_almost\_all\_unconditional}
Over $\Fpt$, the full analytic number theory chain
\[
  \text{Char.\ Cancel.} \;\Rightarrow\; \text{PNT-in-APs}
  \;\Rightarrow\; \FCD \;\Rightarrow\; \SPV
  \;\Rightarrow\; \PSCD \;\Rightarrow\; \text{Almost-All GenMixedMC}
\]
is unconditional. For density-one monic squarefree starting polynomials,
every monic irreducible is reached by some branch of the factor tree.
\end{theorem}

\begin{theorem}[Factor Tree Structure --- Theorem~A]\label{thm:main-tree}
\leancheck{theorem\_A\_factor\_tree}
For any valid starting polynomial and any valid mixed selection:
the walk explores infinitely many distinct monic irreducibles;
the factor pool degree grows to infinity;
the selection is injective (no irreducible appears twice);
the accumulated product is injective;
and the coprimality cascade holds.
\end{theorem}

\begin{theorem}[Stochastic FF-MC --- Theorems~B and~C]\label{thm:main-stochastic}
\leancheck{theorem\_B\_stochastic\_mc}
With any positive probability of non-greedy factor selection at each step,
every monic irreducible is eventually captured. The phase transition is
sharp at $\varepsilon = 0$: deterministic FF-MC is the sole unresolved
boundary.
\end{theorem}

\begin{theorem}[$\Phi_3$ Exclusion Criterion]\label{thm:main-phi3}
\leancheck{autonomous\_map\_landscape}
For $p \equiv 2 \pmod{3}$ with $p \geq 5$, the autonomous map
$f(w) = w(w+1)$ on $\Fp$ has no preimage of~$-1$. The orbit from any
starting value $a \neq -1$ never hits~$-1$.
\end{theorem}

\subsection{Orbit barrier thesis}

Three formally verified dead ends establish the structural boundary of
current techniques:

\begin{enumerate}[label=(\roman*)]
  \item \textbf{Dead End \#127} (Weil insufficient):
    The Weil bound gives population equidistribution but cannot constrain
    orbit dynamics. \leancheck{dead\_end\_127\_weil\_not\_sufficient}
  \item \textbf{Dead End \#129} (FFLM false):
    The large monodromy hypothesis is structurally false---the greedy
    FF-EM construction produces cyclotomic products with abelian Galois
    groups, reducing monodromy arguments to the Weil bound (circular).
    \leancheck{dead\_end\_129\_monodromy\_dead}
  \item \textbf{Dead End \#130} (Generation $\neq$ coverage):
    The steps $\{1,3\}$ generate $\ZZ/4\ZZ$ but the alternating walk
    visits only $\{0,1\}$, missing $\{2,3\}$.
    \leancheck{dead\_end\_130\_generation\_ne\_coverage}
\end{enumerate}

Together, these identify the orbit-specificity barrier as not solvable by
population statistics (Weil, PE, ANT), not solvable by Galois theory
(large monodromy is false), and not solvable by algebraic generation
(subgroup escape $\neq$ dynamical hitting). The sole known dissolution
is stochastic perturbation ($\varepsilon > 0$). The deterministic case
($\varepsilon = 0$) is the open conjecture.

\subsection{Machine verification}

The entire development is formalized in Lean~4 using the Mathlib library
and compiles with zero \texttt{sorry} axioms. The function field module
comprises 17 files totaling 7{,}372 lines of code, with 364 declarations.
We discuss the formalization in Section~\ref{sec:formalization}.

\subsection{Organization}

Section~\ref{sec:background} introduces the algebraic background:
$\Fpt$ as a polynomial ring, monic irreducibles, the necklace identity,
and the autonomous map perspective. Section~\ref{sec:sieve} presents the
unconditional sieve chain culminating in Theorem~\ref{thm:main-sieve}.
Section~\ref{sec:phi3} develops the $\Phi_3$ exclusion criterion
(Theorem~\ref{thm:main-phi3}). Section~\ref{sec:barrier} analyzes the
orbit-specificity barrier and the three dead ends.
Section~\ref{sec:formalization} discusses the Lean formalization.
Section~\ref{sec:open} collects open problems.
